\section*{Capítulo \ref{ch:b2:phanerozoic-earth} --- La Tierra fanerozoica}

\begin{solution}{exo:b2:phanerozoic-earth:1}
(a) París (\ang{5.1}~N) y Nueva York (\ang{3.6}~S); Ciudad del Cabo y Sídney estaban cerca de \ang{60}~S.
(b) Ciudad del Cabo: $61.9 - 33.9 = \ang{28.0}$, $28.0 \times 111.2 = \qty{3114}{km}$; Nueva York:
$40.7 + 3.6 = \ang{44.3}$, \qty{4926}{km}. (c) $3114/300 = \qty{10.4}{km}$ por millón de años,
\qty{1.0}{cm} al año; $4926/300 = \qty{16.4}{km}$ por millón de años, \qty{1.6}{cm} al año.
\end{solution}

\begin{solution}{exo:b2:phanerozoic-earth:2}
(a) 72, 120.1, 50.5 y \qty{135.4}{Myr}; media $(444 - 66)/4 = \qty{94.5}{Myr}$. (b) La
del final del Pérmico (Siberia, a \qty{1.9}{Myr}), la del final del Triásico (CAMP, \qty{0.4}{Myr}) y la del
final del Cretácico (Decán, \qty{0.5}{Myr}). (c) No: los intervalos van de 50 a
\qty{135}{Myr}; una media existe para cualquier conjunto de fechas.
\end{solution}

\begin{solution}{exo:b2:phanerozoic-earth:3}
(a) $\qty{91}{m}/\qty{0.036}{mm.yr^{-1}} = \num{2.5e6}$ años. (b) $3.7/0.036 = 103$ veces. (c)
Lejos de cualquier costa, en aguas claras, poco fango de los ríos llegaba al fondo; la lluvia de
caparazones de algas era el sedimento principal.
\end{solution}

\begin{solution}{exo:b2:phanerozoic-earth:4}
(a) De la arenisca fluvial a la creta, los depósitos se hacen más profundos hacia arriba: una transgresión. La
arena de playa sobre la creta registra una regresión. (b) No: registra una subida del nivel relativo del mar,
que también podrían producir la subsidencia del lugar o una disminución del aporte de sedimento. (c)
Encontrar la misma transgresión, en la misma época, en otros continentes, o eliminar la
subsidencia por \omterm{def:b2:basins-compared:backstripping}{backstripping}.
\end{solution}

\begin{solution}{exo:b2:phanerozoic-earth:5}
(a) $233/97 = \qty{2.4}{m}$ por millón de años, \qty{2.4}{\micro m} al año. (b) $3.7/0.0024 =
1540$: la subida actual es unas 1500 veces más rápida. (c) La curva obtenida por \omterm{def:b2:basins-compared:backstripping}{backstripping}: los sondeos han mostrado
que las cartas de ciclos aciertan en la cronología de los descensos pero exageran su tamaño, en un factor de dos o
más; Nueva Jersey tiene su propio sesgo (la topografía dinámica), de ahí la estimación de $100 \pm
\qty{50}{m}$.
\end{solution}

\begin{solution}{exo:b2:phanerozoic-earth:6}
(a) $\num{e5} \times \num{e9} \times 2900 = \qty{2.9e17}{kg}$ de basalto, y $\num{2.9e17}
\times 0.0075 = \qty{2.2e15}{kg}$ de dióxido de carbono, $\num{2.2e15}/0.04401 = \qty{4.9e16}{mol}$.
(b) $\num{278e-6} \times (44.01/28.97) \times \num{5.1e18} = \qty{2.15e15}{kg}$. (c) Un solo pulso
liberó aproximadamente tanto como contenía todo el aire en 1750. (d) $\num{e17}/\num{2.15e15} = 46.5$: casi cincuenta.
\end{solution}

\begin{solution}{exo:b2:phanerozoic-earth:7}
(a) $2630 \times \frac{\pi}{6} \times (\num{e4})^3 = \qty{1.38e15}{kg}$. (b) $\tfrac12 \times
\num{1.38e15} \times (\num{11.2e3})^2 = \qty{8.7e22}{J}$; a \qty{20}{km.s^{-1}},
\qty{2.76e23}{J}. (c) Un año de \qty{47}{TW} son $\num{47e12} \times \num{3.156e7} =
\qty{1.48e21}{J}$: 59 y 186 años.
\end{solution}

\begin{solution}{exo:b2:phanerozoic-earth:8}
(a) Color oscuro, laminación fina, ausencia de bioturbación, abundante materia orgánica, pirita; a menudo un desplazamiento de los
isótopos del carbono y señales de calentamiento. (b) El evento toarciense con Karoo--Ferrar, el OAE~1a con
Ontong Java. (c) Las fechas tienen incertidumbres de hasta millones de años, y existen muchos eventos y
provincias, así que las coincidencias fortuitas son probables; un vínculo causal requiere que la provincia empiece
justo antes del evento y un trazador del volcanismo, como el mercurio, en los mismos sedimentos.
\end{solution}

\begin{solution}{exo:b2:phanerozoic-earth:9}
(a) $60.9 + 5.1 = \ang{66.0}$, $66.0 \times 111.2 = \qty{7339}{km}$ en \qty{200}{Myr}:
\qty{36.7}{km} por millón de años, \qty{3.7}{cm} al año. (b) Más del doble de los
\qty{1.6}{cm} al año desde \qty{300}{Ma}. (c) Solo se conoce el cambio de latitud; el cambio de
longitud, libre en un marco paleomagnético, y cualquier rodeo se suman a la distancia.
\end{solution}

\begin{solution}{exo:b2:phanerozoic-earth:10}
(a) La fracción de fondo de edad $t$ es $p(t) = \frac2T(1 - t/T)$. Entonces $\int_0^T tp\,\mathrm{d}t =
\frac2T\bigl(\frac{T^2}2 - \frac{T^2}3\bigr) = \frac T3$ y $\int_0^T \sqrt t\,p\,\mathrm{d}t =
\frac2T\bigl(\frac23 T^{3/2} - \frac25 T^{3/2}\bigr) = \frac8{15}\sqrt T$. (b) $T = 3 \times 64.2 =
\qty{192.6}{Myr}$; $2500 + \frac8{15} \times 350 \times 13.878 = \qty{5091}{m}$. (c) $T/2 =
\qty{96.3}{Myr}$: \qty{4332}{m}; el descenso es de \qty{759}{m}, un \qty{2}{\percent} menos que
\qty{774}{m}. (d) Modelo uniforme: $\frac23\sqrt\tau = \frac23\sqrt{2\bar t} = 0.943\sqrt{\bar t}$;
triángulo: $\frac8{15}\sqrt{3\bar t} = 0.924\sqrt{\bar t}$. La respuesta depende de la edad media mucho
más que de la forma.
\end{solution}

\begin{solution}{exo:b2:phanerozoic-earth:11}
(a) En $\Delta s$, solo el factor $1 - f^{-1/2}$ depende de $f$; divídase por su valor en $f = 2$,
$1 - 2^{-1/2} = 0.2929$. (b) $f = (1 - 0.2929\,\Delta s/432)^{-2}$: 1.071 para \qty{50}{m}, 1.239
para \qty{150}{m}; de un 7 a un \qty{24}{\percent} más rápida que hoy. (c) $432 \times (1 -
0.8^{-1/2})/0.2929 = \qty{-174}{m}$: el mar bajaría \qty{174}{m}. (d) El reparto de edades
solo cambia a medida que el fondo antiguo se subduce y se sustituye, lo que lleva un tiempo del orden de $\tau$,
unos \qty{100}{Myr}; el mar lo sigue a lo largo de decenas de millones de años.
\end{solution}

\begin{solution}{exo:b2:phanerozoic-earth:12}
(a) $2^{-7.4/2} = 0.077$: ha desaparecido el \qty{92}{\percent}. (b) $t = 2\log_2(1/0.19) = \qty{4.8}{Myr}$.
(c) $4.8/0.060 = 80$ veces más rápida. (d) Dentro de un piso largo, la extinción de fondo por sí sola elimina
la mayoría de las especies, así que un pulso corto se suma a un total grande y es difícil de separar de él; el
carácter incompleto del registro también desplaza hacia atrás las últimas apariciones.
\end{solution}

\begin{solution}{pb:b2:phanerozoic-earth:1}
\textbf{1.} 2500, 4250, 5304 y \qty{6000}{m}.
\textbf{2.} $\tau = 2 \times 64.2 = \qty{128.4}{Myr}$; $2500 + 350\sqrt{128.4} = \qty{6466}{m}$.
\textbf{3.} $\num{293e6}/(65000 \times 128.4) = \qty{35.1}{km}$ por millón de años,
\qty{35.1}{mm} al año: una velocidad lenta.
\textbf{4.} $\frac1\tau\int_0^\tau(2500 + 350\sqrt t)\,\mathrm{d}t = 2500 + \frac23 \times
350\sqrt\tau = \qty{5144}{m}$.
\textbf{5.} La ley de la raíz cuadrada es demasiado profunda para el fondo de más de unas decenas de millones de años,
que se aplana; las edades reales no están repartidas uniformemente, ya que el fondo joven es el más abundante; y
el sedimento, los \omterm{def:b2:oceanic-islands:seamount}{montes submarinos} y las mesetas elevan el fondo.
\textbf{6.} $\num{293e6} \times 5.144 = \qty{1.51e9}{km^3}$, un \qty{13}{\percent} más que
todo el océano: el modelo es demasiado profundo.
\textbf{7.} $v = \qty{17.6}{km}$ por millón de años; $17.6 \times 128.4 = \qty{2260}{km}$ a
cada lado.
\textbf{8.} $t = 1000/17.6 = \qty{56.8}{Myr}$; $2500 + 350\sqrt{56.8} = \qty{5138}{m}$.
\textbf{9.} En un flanco, con $x = vt$: $\int_0^{v\tau}\bigl(d(\tau) - d(x/v)\bigr)\mathrm{d}x
= vc\int_0^\tau(\sqrt\tau - \sqrt t)\,\mathrm{d}t = \frac13 vc\,\tau^{3/2}$; el doble para ambos.
$\frac23 \times 17.6 \times 0.350 \times 128.4^{3/2} = \frac23 \times 17.6 \times 0.350 \times
1454.9 = \qty{5975}{km^2}$.
\textbf{10.} $65000 \times 5975 = \qty{3.88e8}{km^3}$.
\textbf{11.} $\num{3.88e8}/\num{1.33e9} = 0.29$.
\textbf{12.} Su fondo es en promedio más joven, y por tanto más caliente, menos denso y más alto; la franja de fondo
joven a lo largo del eje es más ancha.
\textbf{13.} $\tau = \qty{64.2}{Myr}$; $2500 + \frac23 \times 350\sqrt{64.2} = \qty{4370}{m}$.
\textbf{14.} $5144 - 4370 = \qty{774}{m}$.
\textbf{15.} $\num{293e6} \times 0.774 = \qty{2.27e8}{km^3}$.
\textbf{16.} $\num{2.27e8}/\num{362e6} = \qty{0.627}{km}$, \qty{627}{m}.
\textbf{17.} $\Delta s = 627 \times (3300 - 1025)/3300 = \qty{432}{m}$.
\textbf{18.} $f = \bigl(1 - 0.2929 \times 100/432\bigr)^{-2} = 1.15$: un \qty{15}{\percent} más rápida.
\textbf{19.} $0.149 \times 148 = \qty{22.1}{\mega km^2}$.
\textbf{20.} $48.1 + 0.32 \times 7.7 = \qty{50.6}{\percent}$, \qty{74.9}{\mega km^2}.
\textbf{21.} Al subir cubre más área, así que el mismo volumen forma una capa más delgada
(\cref{prop:b2:phanerozoic-earth:flooding}).
\textbf{22.} Áreas inundadas a 100, 200, 300 y \qty{400}{m}: 22.1, 41.7, 57.3 y
\qty{71.2}{\mega km^2}. Escalones: $(362 + 11.0) \times 0.1 = \qty{37.3e6}{km^3}$, después 39.4, 41.2
y \qty{42.6e6}{km^3}; acumulados 37.3, 76.7, 117.9 y \qty{160.5e6}{km^3}. El volumen que hay que colocar
es \qty{156.4e6}{km^3}, entre 300 y \qty{400}{m}: $300 + (156.4 - 117.9)/42.6 \times 100
= \qty{390}{m}$.
\textbf{23.} $38.7 + 0.90 \times 9.4 = \qty{47.2}{\percent}$ de la tierra actual.
\textbf{24.} Duplicar la velocidad de expansión haría subir el mar unos \textbf{\qty{390}{m}} y
ahogaría casi la mitad de la tierra actual.
\end{solution}
